Thank you for the insightful questions.

**A1:** Cost-sensitive learning and selective classification are established techniques, but integrating them into a MIL-based histopathology framework with statistical FNR guarantees is non-trivial. Aligning a cost-sensitive objective with Neyman-Pearson calibration and selective abstention while preserving coverage requires careful architectural choices (e.g., a parallel selection head on bag representations) and tailored threshold rules (Cases 1–4). To our knowledge, this specific combination has not been applied to lung adenocarcinoma subtyping with provable FNR control. Furthermore, our ablation results show that no single component alone achieves the observed safety-efficiency trade-offs, justifying the proposed integration.

**A2:** Comparisons with competing models and external validation have been added. See our responses to Reviewer jDt7 (Q12) and Reviewer Pfom (Q2).

**A3:** We will clarify all relevant implementation details in the revised manuscript and our data and code will be publicly released upon acceptance.

**A4:** The prediction head classifies using the task loss, while the selection head learns a sigmoid gating mechanism outputting a coverage probability to decide prediction or abstention. The auxiliary head, used only during training, performs the same classification with different loss weighting — it prevents the selection head from taking easy shortcuts and encourages the representation to retain useful information for both tasks. Jointly optimizing the three heads with the selective risk objective enables the model to abstain on uncertain or high-risk samples, directly addressing the false negative control problem from Phase I.

**A5:** In Phase III, we adopted SelectiveNet with a dedicated selection head rather than simpler confidence-based baselines (e.g., softmax thresholding or temperature scaling). Prior work [12,13] has shown that raw classifier confidence is often miscalibrated, particularly under class imbalance and high-risk medical settings. Moreover, as seen in Phase II (Table 3), even after NP calibration, baseline probabilities yield poor FNR–FPR trade-offs. SelectiveNet jointly optimizes a selection head to minimize selective risk under a coverage constraint, offering a principled reject option (Theorem 2.3). We agree that comparing with softmax thresholding and temperature scaling would strengthen the analysis and will add these comparisons in the revised manuscript.